\chapter{Introduction}

\framac~\cite{userman,sefm12} is a modular analysis framework for the C language
which supports the ACSL specification language~\cite{acsl}. This manual
documents the \eacsl plug-in of \framac, version \eacslversion. Which \eacsl
version you are using is indicated by running \texttt{frama-c
  -e-acsl-version}\optionidx{-}{e-acsl-version}. This plug-in automatically
translates an annotated C program into another program that fails at runtime if
an annotation is violated. If no annotation is violated, the behavior of the new
program is exactly the same than the one of the original program.

Such a translation brings several benefits. First it allows the user to monitor
a \C code, in particular to perform what is usually called ``runtime assertion
checking''~\cite{runtime-assertion-checking}\footnote{In our context, ``runtime
  annotation checking'' would be a better more-general expression.}. This is the
primary goal of \eacsl. Second it allows the combination of \framac and its
existing analyzers, with other analyzer tools for \C like
\pathcrawler~\cite{pathcrawler} that do not natively understand the \acsl
specification language. Third, the possibility to detect invalid annotations
during a concrete execution may be very helpful while writing a correct
specification of a given program, \emph{e.g.} for later program proving.
Finally, an executable specification makes it possible to check assertions that
cannot be verified statically at runtime and thus to establish a link between
monitoring tools and static analysis tools like
\valueplugin~\cite{value}\index{Value} or \wpplugin~\cite{wp}\index{Wp}.

Annotations must be written in the \eacsl specification
language~\cite{eacsl,sac13} which is a subset of \acsl. This plug-in is still in
a preliminary state: some parts of E-ACSL are not yet supported. Which \eacsl
annotations are currently handled by the \eacsl plug-in is documented in a
separated document~\cite{eacsl-implem}.

This manual does \emph{not} explain how to install the plug-in. Please
have a look at file \texttt{INSTALL} of the \eacsl tarball for this
purpose.\index{Installation}
